\IfFileExists{stacks-project.cls}{%
\documentclass{stacks-project}
}{%
\documentclass{amsart}
}

% For dealing with references we use the comment environment
\usepackage{verbatim}
\newenvironment{reference}{\comment}{\endcomment}
%\newenvironment{reference}{}{}
\newenvironment{slogan}{\comment}{\endcomment}
\newenvironment{history}{\comment}{\endcomment}

% For commutative diagrams we use Xy-pic
\usepackage[all]{xy}

% We use 2cell for 2-commutative diagrams.
\xyoption{2cell}
\UseAllTwocells

% We use multicol for the list of chapters between chapters
\usepackage{multicol}

% This is generally recommended for better output
\usepackage{lmodern}
\usepackage[T1]{fontenc}

% For cross-file-references
\usepackage{xr-hyper}

% Package for hypertext links:
\usepackage{hyperref}

% For any local file, say "hello.tex" you want to link to please
% use \externaldocument[hello-]{hello}
\externaldocument[introduction-]{introduction}
\externaldocument[conventions-]{conventions}
\externaldocument[sets-]{sets}
\externaldocument[categories-]{categories}
\externaldocument[topology-]{topology}
\externaldocument[sheaves-]{sheaves}
\externaldocument[sites-]{sites}
\externaldocument[stacks-]{stacks}
\externaldocument[fields-]{fields}
\externaldocument[algebra-]{algebra}
\externaldocument[brauer-]{brauer}
\externaldocument[homology-]{homology}
\externaldocument[derived-]{derived}
\externaldocument[simplicial-]{simplicial}
\externaldocument[more-algebra-]{more-algebra}
\externaldocument[smoothing-]{smoothing}
\externaldocument[modules-]{modules}
\externaldocument[sites-modules-]{sites-modules}
\externaldocument[injectives-]{injectives}
\externaldocument[cohomology-]{cohomology}
\externaldocument[sites-cohomology-]{sites-cohomology}
\externaldocument[dga-]{dga}
\externaldocument[dpa-]{dpa}
\externaldocument[sdga-]{sdga}
\externaldocument[hypercovering-]{hypercovering}
\externaldocument[schemes-]{schemes}
\externaldocument[constructions-]{constructions}
\externaldocument[properties-]{properties}
\externaldocument[morphisms-]{morphisms}
\externaldocument[coherent-]{coherent}
\externaldocument[divisors-]{divisors}
\externaldocument[limits-]{limits}
\externaldocument[varieties-]{varieties}
\externaldocument[topologies-]{topologies}
\externaldocument[descent-]{descent}
\externaldocument[perfect-]{perfect}
\externaldocument[more-morphisms-]{more-morphisms}
\externaldocument[flat-]{flat}
\externaldocument[groupoids-]{groupoids}
\externaldocument[more-groupoids-]{more-groupoids}
\externaldocument[etale-]{etale}
\externaldocument[chow-]{chow}
\externaldocument[intersection-]{intersection}
\externaldocument[pic-]{pic}
\externaldocument[weil-]{weil}
\externaldocument[adequate-]{adequate}
\externaldocument[dualizing-]{dualizing}
\externaldocument[duality-]{duality}
\externaldocument[discriminant-]{discriminant}
\externaldocument[derham-]{derham}
\externaldocument[local-cohomology-]{local-cohomology}
\externaldocument[algebraization-]{algebraization}
\externaldocument[curves-]{curves}
\externaldocument[resolve-]{resolve}
\externaldocument[models-]{models}
\externaldocument[functors-]{functors}
\externaldocument[equiv-]{equiv}
\externaldocument[pione-]{pione}
\externaldocument[etale-cohomology-]{etale-cohomology}
\externaldocument[proetale-]{proetale}
\externaldocument[relative-cycles-]{relative-cycles}
\externaldocument[more-etale-]{more-etale}
\externaldocument[trace-]{trace}
\externaldocument[crystalline-]{crystalline}
\externaldocument[spaces-]{spaces}
\externaldocument[spaces-properties-]{spaces-properties}
\externaldocument[spaces-morphisms-]{spaces-morphisms}
\externaldocument[decent-spaces-]{decent-spaces}
\externaldocument[spaces-cohomology-]{spaces-cohomology}
\externaldocument[spaces-limits-]{spaces-limits}
\externaldocument[spaces-divisors-]{spaces-divisors}
\externaldocument[spaces-over-fields-]{spaces-over-fields}
\externaldocument[spaces-topologies-]{spaces-topologies}
\externaldocument[spaces-descent-]{spaces-descent}
\externaldocument[spaces-perfect-]{spaces-perfect}
\externaldocument[spaces-more-morphisms-]{spaces-more-morphisms}
\externaldocument[spaces-flat-]{spaces-flat}
\externaldocument[spaces-groupoids-]{spaces-groupoids}
\externaldocument[spaces-more-groupoids-]{spaces-more-groupoids}
\externaldocument[bootstrap-]{bootstrap}
\externaldocument[spaces-pushouts-]{spaces-pushouts}
\externaldocument[spaces-chow-]{spaces-chow}
\externaldocument[groupoids-quotients-]{groupoids-quotients}
\externaldocument[spaces-more-cohomology-]{spaces-more-cohomology}
\externaldocument[spaces-simplicial-]{spaces-simplicial}
\externaldocument[spaces-duality-]{spaces-duality}
\externaldocument[formal-spaces-]{formal-spaces}
\externaldocument[restricted-]{restricted}
\externaldocument[spaces-resolve-]{spaces-resolve}
\externaldocument[formal-defos-]{formal-defos}
\externaldocument[defos-]{defos}
\externaldocument[cotangent-]{cotangent}
\externaldocument[examples-defos-]{examples-defos}
\externaldocument[algebraic-]{algebraic}
\externaldocument[examples-stacks-]{examples-stacks}
\externaldocument[stacks-sheaves-]{stacks-sheaves}
\externaldocument[criteria-]{criteria}
\externaldocument[artin-]{artin}
\externaldocument[quot-]{quot}
\externaldocument[stacks-properties-]{stacks-properties}
\externaldocument[stacks-morphisms-]{stacks-morphisms}
\externaldocument[stacks-limits-]{stacks-limits}
\externaldocument[stacks-cohomology-]{stacks-cohomology}
\externaldocument[stacks-perfect-]{stacks-perfect}
\externaldocument[stacks-introduction-]{stacks-introduction}
\externaldocument[stacks-more-morphisms-]{stacks-more-morphisms}
\externaldocument[stacks-geometry-]{stacks-geometry}
\externaldocument[moduli-]{moduli}
\externaldocument[moduli-curves-]{moduli-curves}
\externaldocument[examples-]{examples}
\externaldocument[exercises-]{exercises}
\externaldocument[guide-]{guide}
\externaldocument[desirables-]{desirables}
\externaldocument[coding-]{coding}
\externaldocument[obsolete-]{obsolete}
\externaldocument[fdl-]{fdl}
\externaldocument[index-]{index}

% Theorem environments.
%
\theoremstyle{plain}
\newtheorem{theorem}[subsection]{Theorem}
\newtheorem{proposition}[subsection]{Proposition}
\newtheorem{lemma}[subsection]{Lemma}

\theoremstyle{definition}
\newtheorem{definition}[subsection]{Definition}
\newtheorem{example}[subsection]{Example}
\newtheorem{exercise}[subsection]{Exercise}
\newtheorem{situation}[subsection]{Situation}

\theoremstyle{remark}
\newtheorem{remark}[subsection]{Remark}
\newtheorem{remarks}[subsection]{Remarks}

\numberwithin{equation}{subsection}

% Macros
%
\def\lim{\mathop{\mathrm{lim}}\nolimits}
\def\colim{\mathop{\mathrm{colim}}\nolimits}
\def\Spec{\mathop{\mathrm{Spec}}}
\def\Hom{\mathop{\mathrm{Hom}}\nolimits}
\def\Ext{\mathop{\mathrm{Ext}}\nolimits}
\def\SheafHom{\mathop{\mathcal{H}\!\mathit{om}}\nolimits}
\def\SheafExt{\mathop{\mathcal{E}\!\mathit{xt}}\nolimits}
\def\Sch{\mathit{Sch}}
\def\Mor{\mathop{\mathrm{Mor}}\nolimits}
\def\Ob{\mathop{\mathrm{Ob}}\nolimits}
\def\Sh{\mathop{\mathit{Sh}}\nolimits}
\def\NL{\mathop{N\!L}\nolimits}
\def\CH{\mathop{\mathrm{CH}}\nolimits}
\def\proetale{{pro\text{-}\acute{e}tale}}
\def\etale{{\acute{e}tale}}
\def\QCoh{\mathit{QCoh}}
\def\Ker{\mathop{\mathrm{Ker}}}
\def\Im{\mathop{\mathrm{Im}}}
\def\Coker{\mathop{\mathrm{Coker}}}
\def\Coim{\mathop{\mathrm{Coim}}}

% Boxtimes
%
\DeclareMathSymbol{\boxtimes}{\mathbin}{AMSa}{"02}

%
% Macros for moduli stacks/spaces
%
\def\QCohstack{\mathcal{QC}\!\mathit{oh}}
\def\Cohstack{\mathcal{C}\!\mathit{oh}}
\def\Spacesstack{\mathcal{S}\!\mathit{paces}}
\def\Quotfunctor{\mathrm{Quot}}
\def\Hilbfunctor{\mathrm{Hilb}}
\def\Curvesstack{\mathcal{C}\!\mathit{urves}}
\def\Polarizedstack{\mathcal{P}\!\mathit{olarized}}
\def\Complexesstack{\mathcal{C}\!\mathit{omplexes}}
% \Pic is the operator that assigns to X its picard group, usage \Pic(X)
% \Picardstack_{X/B} denotes the Picard stack of X over B
% \Picardfunctor_{X/B} denotes the Picard functor of X over B
\def\Pic{\mathop{\mathrm{Pic}}\nolimits}
\def\Picardstack{\mathcal{P}\!\mathit{ic}}
\def\Picardfunctor{\mathrm{Pic}}
\def\Deformationcategory{\mathcal{D}\!\mathit{ef}}
\usepackage{cleveref}
\crefname{originallemma}{original lemma}{original lemmas}
\Crefname{originallemma}{Original lemma}{Original lemmas}
\makeatletter
\newlabel{spaces-descent-proposition-fpqc-descent-quasi-coherent}{{04W8}{source}{}{sourcepointer-1}{}}
\newlabel{spaces-descent-proposition-fpqc-descent-quasi-coherent@cref}{{[originallemma][0][]04W8}{[1][0][]source}{}{}{}}
\makeatother
\crefname{originallemma}{original lemma}{original lemmas}
\Crefname{originallemma}{Original lemma}{Original lemmas}
\makeatletter
\newlabel{sites-lemma-mapping-property-glue}{{04TS}{source}{}{sourcepointer-2}{}}
\newlabel{sites-lemma-mapping-property-glue@cref}{{[originallemma][0][]04TS}{[1][0][]source}{}{}{}}
\makeatother
\crefname{originallemma}{original lemma}{original lemmas}
\Crefname{originallemma}{Original lemma}{Original lemmas}
\makeatletter
\newlabel{etale-cohomology-lemma-review-quasi-coherent}{{0DEW}{source}{}{sourcepointer-3}{}}
\newlabel{etale-cohomology-lemma-review-quasi-coherent@cref}{{[originallemma][0][]0DEW}{[1][0][]source}{}{}{}}
\makeatother
\crefname{originallemma}{original lemma}{original lemmas}
\Crefname{originallemma}{Original lemma}{Original lemmas}
\makeatletter
\newlabel{sites-modules-section-local}{{03DK}{source}{}{sourcepointer-4}{}}
\newlabel{sites-modules-section-local@cref}{{[originallemma][0][]03DK}{[1][0][]source}{}{}{}}
\makeatother
\crefname{originallemma}{original lemma}{original lemmas}
\Crefname{originallemma}{Original lemma}{Original lemmas}
\makeatletter
\newlabel{sites-definition-special-cocontinuous-functor}{{03CG}{source}{}{sourcepointer-5}{}}
\newlabel{sites-definition-special-cocontinuous-functor@cref}{{[originallemma][0][]03CG}{[1][0][]source}{}{}{}}
\makeatother
\crefname{originallemma}{original lemma}{original lemmas}
\Crefname{originallemma}{Original lemma}{Original lemmas}
\makeatletter
\newlabel{sites-lemma-equivalence}{{03A0}{source}{}{sourcepointer-6}{}}
\newlabel{sites-lemma-equivalence@cref}{{[originallemma][0][]03A0}{[1][0][]source}{}{}{}}
\makeatother
\title{Review of quasi-coherent modules on the etale and fppf sites of an algebraic space}
\author{The Stacks Project Authors}
\date{stacks-project:spaces-more-cohomology.tex:local-snapshot-20261012; GNU FDL 1.2 or later}
\begin{document}
\maketitle
\noindent\textit{Editorial note.} The counted statement is the author's own lemma environment labelled lemma-review-quasi-coherent (source lines 1519-1546); the proof environment beginning on the line after the following blank line (source line 1548) is its own proof and closes at source line 1596.\par\medskip
\noindent\textit{Editorial note.} The proof invokes one result of the same chapter by label (retained as context with the generator's transitive expansion) and four cross-chapter results (Descent on Spaces, Sites twice, Etale Cohomology, Modules on Sites), each declared as a tagged source pointer.\par\medskip
\noindent\textit{Editorial note.} The author closes with 'Some minor details omitted.'. This is the author's own closing remark on the standard transport of quasi-coherence along the equivalences; the counted proof itself is a complete 49-line reduction (fpqc descent, locality over an etale covering by schemes, and the scheme case), so the remark is recorded as written rather than treated as an omitted proof.\par\medskip
\noindent\textit{Editorial note.} The retained preamble is the chapter's own (sources/stacks-spaces-more-cohomology/preamble.tex, taken with its real line count), a verbatim excerpt of the native Xy-pic setup and theorem declarations; no macro was invented or redefined.\par\medskip
\noindent\textit{Editorial scope.} One original result of the Stacks Project's chapter Cohomology of Algebraic Spaces, counted once: the seven-part review of quasi-coherent modules on the etale and fppf sites of an algebraic space.\par\medskip
\noindent\textit{Allowed standard prerequisites (editorial).}\begin{itemize}
\item fpqc descent of quasi-coherent modules. Quasi-coherent modules on algebraic spaces satisfy fpqc descent: compatible descent data descend uniquely (Descent on Spaces, declared as a tagged source pointer).. It is the whole of part (1) at source lines 1549-1564.
\item Sheaves as descent data for a covering of the final object. If a covering of the final object is given, the category of sheaves is equivalent to the category of descent data for that covering (Sites, declared as a tagged source pointer).. It makes parts (2)-(7) local on X at source lines 1568-1576.
\item Equivalence of ringed topoi for schemes versus spaces. The embeddings of the etale (and fppf) sites of schemes into those of spaces determine equivalences of ringed topoi (the author's lemma-compare-cohomology-other-topologies, retained as context).. It transfers the scheme case to the general case at source lines 1583-1588.
\item The scheme case and intrinsicness of quasi-coherence. The corresponding seven-part statement for schemes (Etale Cohomology, declared as a tagged source pointer) and the fact that quasi-coherence is an intrinsic property of modules, so it transports along the equivalence (Modules on Sites, declared as a tagged source pointer).. They finish the proof at source lines 1589-1594.
\end{itemize}\medskip
\ifdefined\hypertarget\hypertarget{sourcepointer-1}{}\fi
\noindent\textit{Source reference.} Original lemma 04W8: sources/stacks-chow/tags.txt, source lines 6351--6351. Cross-chapter pointer whose target carries the canonical tag 04W8 in the project tag list; declared so the reference resolves inside the assembled package, while every step the counted argument uses remains inside the retained spans.\par\medskip
\ifdefined\hypertarget\hypertarget{sourcepointer-2}{}\fi
\noindent\textit{Source reference.} Original lemma 04TS: sources/stacks-chow/tags.txt, source lines 6265--6265. Cross-chapter pointer whose target carries the canonical tag 04TS in the project tag list; declared so the reference resolves inside the assembled package, while every step the counted argument uses remains inside the retained spans.\par\medskip
\ifdefined\hypertarget\hypertarget{sourcepointer-3}{}\fi
\noindent\textit{Source reference.} Original lemma 0DEW: sources/stacks-chow/tags.txt, source lines 16816--16816. Cross-chapter pointer whose target carries the canonical tag 0DEW in the project tag list; declared so the reference resolves inside the assembled package, while every step the counted argument uses remains inside the retained spans.\par\medskip
\ifdefined\hypertarget\hypertarget{sourcepointer-4}{}\fi
\noindent\textit{Source reference.} Original lemma 03DK: sources/stacks-chow/tags.txt, source lines 4423--4423. Cross-chapter pointer whose target carries the canonical tag 03DK in the project tag list; declared so the reference resolves inside the assembled package, while every step the counted argument uses remains inside the retained spans.\par\medskip
\ifdefined\hypertarget\hypertarget{sourcepointer-5}{}\fi
\noindent\textit{Source reference.} Original lemma 03CG: sources/stacks-chow/tags.txt, source lines 4383--4383. Cross-chapter pointer whose target carries the canonical tag 03CG in the project tag list; declared so the reference resolves inside the assembled package, while every step the counted argument uses remains inside the retained spans.\par\medskip
\ifdefined\hypertarget\hypertarget{sourcepointer-6}{}\fi
\noindent\textit{Source reference.} Original lemma 03A0: sources/stacks-chow/tags.txt, source lines 4291--4291. Cross-chapter pointer whose target carries the canonical tag 03A0 in the project tag list; declared so the reference resolves inside the assembled package, while every step the counted argument uses remains inside the retained spans.\par\medskip
\begin{lemma}
\label{lemma-compare-cohomology-other-topologies}
Let $S$ be a scheme. Let $\tau \in \{\etale, fppf, ph\}$ (add more here).
The inclusion functor
$$
(\Sch/S)_\tau \longrightarrow (\textit{Spaces}/S)_\tau
$$
is a special cocontinuous functor
(Sites, Definition \ref{sites-definition-special-cocontinuous-functor})
and hence identifies topoi.
\end{lemma}

\begin{proof}
The conditions of Sites, Lemma \ref{sites-lemma-equivalence}
are immediately verified as our functor is fully faithful
and as every algebraic space has an \'etale covering by schemes.
\end{proof}

\begin{lemma}
\label{lemma-review-quasi-coherent}
Let $S$ be a scheme. Let $X$ be an algebraic space over $S$.
Let $\mathcal{F}$ be a quasi-coherent $\mathcal{O}_X$-module.
\begin{enumerate}
\item The rule
$$
\mathcal{F}^a : (\textit{Spaces}/X)_\etale \longrightarrow \textit{Ab},\quad
(f : Y \to X) \longmapsto \Gamma(Y, f^*\mathcal{F})
$$
satisfies the sheaf condition for fpqc and a fortiori
fppf and \'etale coverings,
\item $\mathcal{F}^a = \pi_X^*\mathcal{F}$ on $(\textit{Spaces}/X)_\etale$,
\item $\mathcal{F}^a = a_X^*\mathcal{F}$ on $(\textit{Spaces}/X)_{fppf}$,
\item the rule $\mathcal{F} \mapsto \mathcal{F}^a$ defines
an equivalence between quasi-coherent $\mathcal{O}_X$-modules
and quasi-coherent modules on
$((\textit{Spaces}/X)_\etale, \mathcal{O})$,
\item the rule $\mathcal{F} \mapsto \mathcal{F}^a$ defines
an equivalence between quasi-coherent $\mathcal{O}_X$-modules
and quasi-coherent modules on
$((\textit{Spaces}/X)_{fppf}, \mathcal{O})$,
\item we have $\epsilon_{X, *}a_X^*\mathcal{F} = \pi_X^*\mathcal{F}$
and $a_{X, *}a_X^*\mathcal{F} = \mathcal{F}$,
\item we have $R^i\epsilon_{X, *}(a_X^*\mathcal{F}) = 0$
and $R^ia_{X, *}(a_X^*\mathcal{F}) = 0$ for $i > 0$.
\end{enumerate}
\end{lemma}

\begin{proof}
Part (1) is a consequence of fppf descent of quasi-coherent modules.
Namely, suppose that $\{f_i : U_i \to U\}$ is an fpqc covering
in $(\textit{Spaces}/X)_\etale$. Denote $g : U \to X$ the structure
morphism. Suppose that
we have a family of sections $s_i \in \Gamma(U_i , f_i^*g^*\mathcal{F})$
such that $s_i|_{U_i \times_U U_j} = s_j|_{U_i \times_U U_j}$.
We have to find the correspond section $s \in \Gamma(U, g^*\mathcal{F})$.
We can reinterpret the $s_i$ as a family of maps
$\varphi_i : f_i^*\mathcal{O}_U = \mathcal{O}_{U_i} \to f_i^*g^*\mathcal{F}$
compatible with the canonical descent data associated to the
quasi-coherent sheaves $\mathcal{O}_U$ and $g^*\mathcal{F}$ on $U$.
Hence by Descent on Spaces, Proposition
\ref{spaces-descent-proposition-fpqc-descent-quasi-coherent}
we see that we may (uniquely) descend
these to a map $\mathcal{O}_U \to g^*\mathcal{F}$ which gives
us our section $s$.

\medskip\noindent
We will deduce (2) -- (7) from the corresponding statement for schemes.
Choose an \'etale covering $\{X_i \to X\}_{i \in I}$
where each $X_i$ is a scheme. Observe that $X_i \times_X X_j$
is a scheme too. This covering induces a covering of
the final object in each of the three sites
$(\textit{Spaces}/X)_{fppf}$, $(\textit{Spaces}/X)_\etale$, and $X_\etale$.
Hence we see that the category of sheaves on these sites
are equivalent to descent data for these coverings, see
Sites, Lemma \ref{sites-lemma-mapping-property-glue}.
Parts (2), (3) are local (because we have the glueing statement).
Being quasi-coherent is a local property, hence parts
(4), (5) are local. Clearly (6) and (7) are local.
It follows that it suffices to prove parts (2) -- (7)
of the lemma when $X$ is a scheme.

\medskip\noindent
Assume $X$ is a scheme. The embeddings
$(\Sch/X)_\etale \subset (\textit{Spaces}/X)_\etale$ and
$(\Sch/X)_{fppf} \subset (\textit{Spaces}/X)_{fppf}$
determine equivalences of ringed topoi by
Lemma \ref{lemma-compare-cohomology-other-topologies}.
We conclude that (2) -- (7) follows from the case of schemes.
\'Etale Cohomology, Lemma
\ref{etale-cohomology-lemma-review-quasi-coherent}.
To transport the property of being quasi-coherent via
this equivalence use that being quasi-coherent is an
intrinsic property of modules as explained in
Modules on Sites, Section \ref{sites-modules-section-local}.
Some minor details omitted.
\end{proof}

\end{document}
